% ============================================================================
% Fichier  : partie5/P05_C16_mesures_techniques.tex
% Titre    : Mesures Techniques
% Auteur   : MINKA MI NGUIDJOÏ Thierry Emmanuel
% ============================================================================
% Sources : ISO/IEC 27002:2022 (thème technologique) ; ANSSI. Fictif : Eséka.
% ============================================================================

\chapter{Mesures Techniques}
\label{chap:technique}
\chaptermark{Mesures techniques}

\epigraph{%
    « Une bonne mesure technique applique une décision de sécurité~;\\
    elle ne la remplace pas. »%
}{Principe du cours}

\niveauChapitre{Intermédiaire}{Le Chapitre~\ref{chap:orga}~; notions de réseaux.}

\begin{boxobjectifs}
\begin{itemize}
    \item Distinguer les modèles de contrôle d'accès (\sigle{DAC}, \sigle{MAC},
          \sigle{RBAC}).
    \item Expliquer l'authentification forte (\sigle{MFA}).
    \item Situer pare-feu, \sigle{VPN} et segmentation réseau.
    \item Comprendre le rôle du chiffrement et de la journalisation.
    \item Relier chaque mesure technique à un critère \sigle{DIC}.
\end{itemize}
\end{boxobjectifs}

\bigskip

% ============================================================================
\section*{Question directrice : quelles mesures techniques pour quels risques ?}
% ============================================================================

Les mesures techniques sont nombreuses~: le danger est de les empiler sans
logique. Ce chapitre les présente en les rattachant aux risques et aux critères
de sécurité de la mairie.

\medskip
\textbf{Plan du chapitre.}
\begin{enumerate}
    \item Le contrôle d'accès (section~\ref{sec:c16-acces}).
    \item L'authentification forte (section~\ref{sec:c16-mfa}).
    \item La sécurité réseau (section~\ref{sec:c16-reseau}).
    \item Chiffrement et journalisation (section~\ref{sec:c16-crypto}).
    \item Rattacher aux critères (section~\ref{sec:c16-dic}).
\end{enumerate}

\separateur

% ============================================================================
\section{Le Contrôle d'Accès}
\label{sec:c16-acces}
% ============================================================================

\begin{tableaucours}{L{2.2cm}L{4.4cm}L{5.6cm}}{%
    \tableauHeader{Modèle} &
    \tableauHeader{Principe} &
    \tableauHeader{Exemple}
}
DAC & Fondé sur l'identité, propriétaire & Permissions d'un fichier \\
MAC & Fondé sur des étiquettes & Politique centralisée (SELinux) \\
RBAC & Fondé sur les rôles & Un agent d'état civil voit l'état civil \\
\end{tableaucours}
\vspace{-0.8em}
\begin{center}{\small\itshape Tableau~16.1. Les modèles de contrôle d'accès}\end{center}

Le \termecle{moindre privilège} commande~: chacun n'a que les droits
nécessaires à sa tâche. À la mairie, le modèle \sigle{RBAC} convient~: on donne
des droits par rôle (état civil, taxes), non par personne.

% ============================================================================
\section{L'Authentification Forte}
\label{sec:c16-mfa}
% ============================================================================

\begin{boxdef}
\termecle{Authentification multifacteur} (\sigle{MFA})~: vérification de
l'identité par au moins deux facteurs parmi~: ce que l'on sait (mot de passe),
ce que l'on possède (jeton, téléphone), ce que l'on est (biométrie).
\end{boxdef}

La \sigle{MFA} réduit fortement le risque lié aux mots de passe volés ou
faibles. Pour les accès sensibles de la mairie (administration du serveur
d'état civil), elle devrait être la règle.

% ============================================================================
\section{La Sécurité Réseau}
\label{sec:c16-reseau}
% ============================================================================

\begin{tableaucours}{L{3.4cm}L{8.6cm}}{%
    \tableauHeader{Mesure} &
    \tableauHeader{Rôle}
}
Pare-feu & Filtrer les flux, bloquer par défaut \\
Segmentation & Isoler les réseaux (postes publics / serveurs) \\
VPN & Chiffrer les accès distants \\
IDS/IPS & Détecter ou bloquer les intrusions \\
\end{tableaucours}
\vspace{-0.8em}
\begin{center}{\small\itshape Tableau~16.2. Mesures de sécurité réseau}\end{center}

La \termecle{segmentation} aurait limité l'incident d'Eséka~: si les postes du
public avaient été isolés du serveur d'état civil, la clé n'aurait pas tout
atteint. Cloisonner, c'est contenir.

% ============================================================================
\section{Chiffrement et Journalisation}
\label{sec:c16-crypto}
% ============================================================================

\begin{boxdef}
\termecle{Chiffrement}~: rendre une information inintelligible sans la clé,
pour protéger sa \emph{confidentialité}. \termecle{Journalisation}~:
enregistrer les événements du système, pour la \emph{traçabilité} et la
détection.
\end{boxdef}

Chiffrer le fichier des contribuables protège sa confidentialité même en cas de
vol du support. Journaliser les accès à l'état civil permet de savoir qui a
consulté quoi, et soutient l'enquête en cas d'incident (lien avec le cours de
forensique, \sigle{LPI342}).

\begin{boxmetier}{DSI}
\textbf{Déployer ce que la politique décide.} Le \sigle{DSI} met en \oe uvre,
mais ne décide pas seul~: il déploie les mesures que l'appréciation des risques
a justifiées et que la politique a actées. La technique sert la décision, pas
l'inverse.
\end{boxmetier}

% ============================================================================
\section{Rattacher aux Critères}
\label{sec:c16-dic}
% ============================================================================

\begin{boxprocedure}
\textbf{Constatation.} La mairie déploie~: \sigle{RBAC}, \sigle{MFA},
segmentation, chiffrement, journalisation, sauvegarde.

\medskip
\textbf{Analyse.} Chaque mesure sert un critère~: contrôle d'accès et
chiffrement servent la \emph{confidentialité}~; sauvegarde et segmentation, la
\emph{disponibilité}~; journalisation, la \emph{traçabilité}~; signature et
contrôle d'intégrité, l'\emph{intégrité}. Rattacher chaque mesure à un critère
évite d'en oublier et d'en empiler d'inutiles.
\end{boxprocedure}

\separateur

\begin{boxresume}
\textbf{Ce qu'il faut retenir de ce chapitre~:}
\begin{itemize}
    \item \textbf{DAC, MAC, RBAC}~: trois modèles de contrôle d'accès~; le
          \textbf{moindre privilège} commande.
    \item La \textbf{MFA} réduit le risque des mots de passe volés.
    \item \textbf{Pare-feu, segmentation, VPN} cloisonnent et protègent le
          réseau.
    \item \textbf{Chiffrement} (confidentialité) et \textbf{journalisation}
          (traçabilité) sont essentiels.
    \item Chaque mesure technique se rattache à un \textbf{critère DIC}.
\end{itemize}
\end{boxresume}

\begin{boxinfo}
\textbf{Points de contrôle de l'élève.} Je sais~: distinguer DAC, MAC, RBAC~;
expliquer la MFA~; situer les mesures réseau~; décrire chiffrement et
journalisation~; rattacher une mesure à un critère.
\end{boxinfo}

\bigskip

\begin{boxexo}[title={\ding{45}\quad Exercice~16.1. Les modèles \niveaubase}]
Distinguez DAC, MAC et RBAC en une phrase chacun.
\end{boxexo}

\begin{boxexo}[title={\ding{45}\quad Exercice~16.2. MFA \niveaubase}]
Qu'est-ce que la MFA, et quels sont les trois types de facteurs~?
\end{boxexo}

\begin{boxexo}[title={\ding{45}\quad Exercice~16.3. Rôle d'une mesure \niveaubase}]
À quoi sert la segmentation réseau~?
\end{boxexo}

\begin{boxexo}[title={\ding{45}\quad Exercice~16.4. Moindre privilège \niveaumoyen}]
Appliquez le moindre privilège aux agents d'état civil et des taxes.
\end{boxexo}

\begin{boxexo}[title={\ding{45}\quad Exercice~16.5. Segmentation \niveaumoyen}]
Comment la segmentation aurait-elle limité l'incident de la clé \texttt{USB}~?
\end{boxexo}

\begin{boxexo}[title={\ding{45}\quad Exercice~16.6. Chiffrer quoi \niveaumoyen}]
Quelles informations de la mairie chiffreriez-vous en priorité, et pourquoi~?
\end{boxexo}

\begin{boxexo}[title={\ding{45}\quad Exercice~16.7. Journaliser \niveaumoyen}]
Quels événements journaliser sur le serveur d'état civil~?
\end{boxexo}

\begin{boxexo}[title={\ding{45}\quad Exercice~16.8. Mesure et critère \niveaumoyen}]
Reliez chaque mesure à un critère DIC~: chiffrement, sauvegarde,
journalisation, RBAC.
\end{boxexo}

\begin{boxexo}[title={\ding{45}\quad Exercice~16.9. Architecture \niveauexpert}]
Proposez une architecture réseau segmentée pour la mairie d'Eséka (public,
agents, serveurs).
\end{boxexo}

\begin{boxexo}[title={\ding{45}\quad Exercice~16.10. Plan technique \niveauexpert}]
Dressez le plan des mesures techniques de la mairie, chacune rattachée à un
risque et à un critère.
\end{boxexo}

\bigskip

\begin{boxinfo}
\textbf{Pour aller plus loin.}
\begin{itemize}
    \item ISO/IEC 27002, thème technologique~\cite{iso27002}~; Guide d'hygiène
          ANSSI~\cite{anssihygiene2017}.
\end{itemize}
\end{boxinfo}

\begin{boxinfo}
\textbf{Transition.} Reste à se préparer au pire~: le Chapitre~\ref{chap:continuite}
traite la continuité d'activité, \sigle{PCA} et \sigle{PRA}.
\end{boxinfo}
